\section{Question 6}

\subsection{Question 6.1}

\subsubsection*{a.}

The True Positive Rate (TPR) is defined as follows:

\[TPR = \frac{TP}{TP + FN},\]

where $TP$ is the True Positives and $FN$ is the False Negatives. It is also known as
the Recall or Sensitivity. 

The False Positive Rate (FPR) is defined as follows:

\[FPR = \frac{FP}{FP + TN},\]

where $FP$ is the False Positives and $TN$ is the True Negatives.

The Receiver Operating Characteristic (ROC) curve, plots the TPR against the FPR across
varying classification thresholds. Moving along the curve from left to right represents decreasing the threshold from
$+\infty$ to $-\infty$. Hence, a low threshold catches more churners but also results in more
false alarms. A higher threshold reduces false alarms but ends up missing more actual churners.

\subsubsection*{b.}

An AUC of 0.5 represents a random classifier. Hence it is no better than simply randomly guessing.
In terms of the churn detection task it represents a classifier that has no real ability to 
discriminate churners from non-churners.

Whereas, an AUC of 1 represents a perfect classifier. In terms of the churn detection task
this would indicate a classifier that is able to perfectly discriminate between non-churners
and churners.

\subsubsection*{c.}

In this particular case, due to the constraint on the amount of retention offers that can be
sent, precision would be a better metric than AUC. AUC is good for evaluating a model's 
overall performance across thresholds but it does not directly tell the company how many of the
selected clients will actually churn. Precision does this by indicating which of the users
will actually churn from the all the predicted churners.

In this case the Precision would be:

\[Precision = \frac{TP}{TP + FP} = \frac{60}{60 + 150} \approx 0.286\]

This would indicate that of the flagged customers, only 28.6\% were actual churners. By optimising
the Precision at a certain Recall level, one can reduce the false alarms and ensure the retention 
offers stay below the threshold.

\subsection{Question 6.2}

Figure \ref{fig:roc} illustrates the ROC curves of the SGD logistic regression, RBF-SVM and
Ridge logistic regression models. The AUC score for each model is also shown.

\begin{figure}[h]
    \centering
    \includegraphics[width=0.75\textwidth]{Images/roc.png}
    \caption{ROC Curves of the three Trained Models}
    \label{fig:roc}
\end{figure}

From Figure \ref{fig:roc}, the Ridge logistic regression model has the best overall performance
across threshold since it has the highest AUC of 0.611. Hence, this would be the model that should
be deployed since it have the best ability out of the three models to rank churners above
non-churners across different classification thresholds. However, the AUC value is still quite
low and hence, this is a weak classifier. Better features, more data and/or more fine-tuning would be required 
before deploying such a model in a real-world churn prediction situation.
